\documentclass[10pt,a4paper]{amsart}
\usepackage[margin=19mm]{geometry}
\usepackage{amsmath,amssymb,mathtools}
\usepackage[scr=rsfs,cal=euler]{mathalfa}
\usepackage{xfrac}
\usepackage[unicode,hidelinks,bookmarks=false]{hyperref}
\newcommand{\ie}{i.e.\ }
\newcommand{\cf}{cf.\ }
\newcommand{\resp}{respectively\ }
\newcommand{\Subsection}{\subsection}
\providecommand{\nobreakdash}{\nobreak\hbox{-}}
\setlength{\emergencystretch}{2em}
\multlinegap0pt
\usepackage{amssymb}
\usepackage{mathtools}
\usepackage{tikz}
\usepackage{xcolor}
\usepackage{float}
\newtheorem{lm}{Lemma}
\newtheorem{thm}{Theorem}
\title{Partitions of a rectangle}
\author{Krystian Gajdzica, Maciej Zakarczemny}
\date{Source: arXiv:2608.08762v1 (CC BY 4.0)}
\begin{document}
\maketitle

\noindent\emph{Source and licence.} Partitions of a rectangle. Version arXiv:2608.08762v1 (CC BY 4.0), distributed by arXiv under the Creative Commons Attribution 4.0 International licence, http://creativecommons.org/licenses/by/4.0/, as recorded in the arXiv licence metadata for this exact version and shown on the abstract page https://arxiv.org/abs/2608.08762v1. Full licence notice: \texttt{licenses/arxiv-2608.08762-CC-BY-4.0.txt}.\par\smallskip

\noindent\textit{Editorial scope.} Counted unit: the article's thm Gc (source0.tex lines 4093--4111). The article's complete original proof of it is delivered with it (source0.tex lines 4113--4200, one \texttt{proof} environment of 88 lines). No mathematics is written, repaired or re-derived: the statement and the proof are the article's own lines, in the article's own notation, and no retained line is substituted or re-flowed.  Retained uncounted as the source-local context the proof needs: lines 3600--3626; lines 3786--3831.  Nothing was omitted from any retained span, so the package carries no omission record.  No retained line contains a citation, so the wrapper carries no bibliography and no bibliography entry.  No retained line contains a figure environment, an image-inclusion command or drawing code, so no image is delivered and the package declares no asset dependency.  Editorial formatting only, in addition: an AMS \texttt{amsart} preamble that restates the article's own declarations and macros used by the retained lines, namely the environments \texttt{lm}, \texttt{thm} on separate counters, together with the standard packages those lines need.  The retained environments print in this document as Lm 1, Theorem 1 in order of appearance, whereas the article prints the same items with the numbering of its own full document.  The complete native source of the article as distributed in the arXiv e-print for this exact version is bundled unchanged as \texttt{sources/207187/source0.tex}.
\begin{lm}\label{Gcomasym}
Let us fix $\Delta>0$ and put
$$
G(x)=
\tfrac{1}{(1-x^3)(1-x^6)}
\left(\prod_{j=1}^\infty\tfrac{1}{1-x^j}\right)
\left(\prod_{j=1}^\infty\tfrac{1}{1-x^{2j}}\right).
$$
As the complex variable $z\to0$ in the region
$
\operatorname{Re}z>0
$
and
$
|\operatorname{Im}z|
\le 
\Delta\operatorname{Re}z,
$
we have that
$$
G(e^{-z})
\sim
\tfrac{1}{18\pi\sqrt2}\,z^{-1}
\exp\left(\tfrac{\pi^2}{4z}\right),
$$
uniformly in that region.
\end{lm}

\begin{thm}\label{thm: BJSM}
Suppose that $B(x)=\sum_{j=0}^\infty b_jx^j$ is a power series with non-negative real coefficients and radius of convergence at least one. Let $\lambda,\gamma>0$ and $\alpha,\beta\in\mathbb{R}$. Suppose that
\begin{align*}
B(e^{-t})
&\sim
\lambda
\left(\log\frac1t\right)^\alpha
t^\beta
\exp\left(\frac{\gamma}{t}\right),
\qquad\text{as }t\to0^+,\\
B(e^{-z})
&=
O\left(
\left(\log\frac1{|z|}\right)^\alpha
|z|^\beta
\exp\left(\frac{\gamma}{|z|}\right)
\right),
\qquad\text{as }z\to0,
\end{align*}
where $z=x+iy$, $x,y\in\mathbb{R}$, $x>0$, and the second estimate holds in every region of the~form $|y|\le \Delta x$ with $\Delta>0$. Then
$$
\sum_{k=0}^n b_k
\sim
\frac{
\lambda\gamma^{\frac{\beta}{2}-\frac{1}{4}}(\log n)^\alpha
}{
2^{\alpha+1}\sqrt{\pi}\,
n^{\frac{\beta}{2}+\frac{1}{4}}
}
\exp\left(2\sqrt{\gamma n}\right),
\qquad\text{as }n\to\infty.
$$
Furthermore, if the sequence $(b_n)_{n\geq0}$ is weakly increasing, then
$$
b_n
\sim
\frac{
\lambda\gamma^{\frac{\beta}{2}+\frac{1}{4}}(\log n)^\alpha
}{
2^{\alpha+1}\sqrt{\pi}\,
n^{\frac{\beta}{2}+\frac{3}{4}}
}
\exp\left(2\sqrt{\gamma n}\right),
\qquad\text{as }n\to\infty.
$$
\end{thm}

\begin{thm}\label{Gc}
Let
\begin{equation}\label{eq:G-coefficients}
G(x)=\frac{1}{(1-x^3)(1-x^6)}
\left(\prod_{n=1}^\infty\frac{1}{1-x^n}\right)
\left(\prod_{n=1}^\infty\frac{1}{1-x^{2n}}\right)
=
\sum_{n=0}^{\infty}b_nx^n.
\end{equation}
Then
$$
b_n
\sim
\tfrac{1}{36\pi^2}
n^{-\frac{1}{4}}
\exp\left(\pi\sqrt n\right),
\hspace{0.5cm}\text{as }n\to\infty.
$$
\end{thm}

\begin{proof}
We apply Theorem~\ref{thm: BJSM} to the generating function $G(x)$. First, we can write\\ $G(x)=\frac{1}{1-x}H(x),$
where
$$
H(x)=\frac{1}{(1-x^3)(1-x^6)}
\left(\prod_{n=2}^\infty\frac{1}{1-x^n}\right)
\left(\prod_{n=1}^\infty\frac{1}{1-x^{2n}}\right).
$$
The coefficients of both $G(x)$ and $H(x)$ are non-negative. Moreover, if
$H(x)=\sum_{n=0}^{\infty}h_nx^n,$
then
$b_n=\sum_{j=0}^{n}h_j.
$
In particular, this implies that the sequence $(b_n)_{n\ge0}$ is weakly increasing.

For every $\Delta>0$, Lemma~\ref{Gcomasym} guarantees that
$$
G(e^{-z})
\sim
\tfrac{1}{18\pi\sqrt2}\,
z^{-1}
\exp\left(\tfrac{\pi^2}{4z}\right)
$$
uniformly as $z\to0$ in the region
$
\operatorname{Re}z>0
$
and
$
|\operatorname{Im}z|\le \Delta\operatorname{Re}z.
$
This uniform asymptotic also gives the required complex estimate. Indeed, one has
$$
G(e^{-z})
=
O\!\left(
|z|^{-1}
\left|
\exp\left(\tfrac{\pi^2}{4z}\right)
\right|
\right)
$$
for sufficiently small $z$ in the region. Since
$
\left|
\exp\left(\tfrac{\pi^2}{4z}\right)
\right|
=
\exp\left(
\frac{\pi^2}{4}
\operatorname{Re}\frac{1}{z}
\right)
$
and
$
\operatorname{Re}\frac{1}{z}
\le 
\frac{1}{|z|},
$
we obtain that
$$
G(e^{-z})=
O\left(
|z|^{-1}
\exp\left(\tfrac{\pi^2}{4|z|}\right)\right).
$$
Therefore, the assumptions of Theorem \ref{thm: BJSM} are satisfied with
\begin{align*}
    \lambda=\tfrac{1}{18\pi\sqrt2},\hspace{0.5cm} \alpha=0,\hspace{0.5cm}\beta=-1\hspace{0.5cm}\text{and}\hspace{0.5cm}\gamma=\tfrac{\pi^2}{4}.
\end{align*}
Hence, we have that
$$
b_n
\sim
\tfrac{\lambda}{2\sqrt\pi}
\gamma^{\tfrac{\beta}{2}+\tfrac{1}{4}}
n^{-\tfrac{\beta}{2}-\tfrac{3}{4}}
\exp\left(2\sqrt{\gamma n}\right).
$$
After substituting, one can conclude that
$$
b_n
\sim
\tfrac{1}{36\pi^2}
n^{-\frac{1}{4}}
\exp\left(\pi\sqrt n\right).
$$
\end{proof}

\end{document}
